\section{Evaluation contamination and uncertainty}\label{sec:C16:evaluation}
A retrieval evaluation should specify the unit of relevance and the unit of success. Passage relevance can be judged independently of answer generation, while task success includes retrieval, reasoning, and execution. Report both when possible. An oracle evidence condition identifies a useful ceiling, but it must be labeled as diagnostic because a real system does not know the decisive evidence in advance.

ReFind supplies a concrete counterweight to assuming that more elaborate memory structure is necessarily better. It retains raw chat turns, uses lexical indexing, and gives an iterative agent controls for session-aware rank fusion, neighborhood expansion, time filtering, and previously inspected sessions; a separate stage answers from saved evidence~\cite{cumb2026refindv2}. The absence of generated semantic structure does not mean the absence of organization: session boundaries and timestamps remain essential. Its reused leaderboard baselines should be distinguished from its matched control experiments. The scientific comparison is thus not ``raw versus structured'' in isolation, but offline abstraction versus query-time search under specified access controls, backbone, final evidence, and complete cost. This baseline is especially relevant before attributing gains to a graph or to expensive LLM-based index construction.

Use retrieval ablations that preserve final context length. Comparing a graph method returning many rich records with a flat baseline returning fewer tokens confounds structure and evidence volume. Similarly, comparing an active retriever with a one-shot baseline under equal retrieval-call limits can conceal additional model computation. Match either complete resource budgets or report a Pareto surface over success, token use, tool interactions, latency, and construction cost. One scalar rank rarely captures all relevant tradeoffs.

Contamination can enter through the index, training data, task generator, and reranker labels. A trajectory index containing near-duplicate evaluation solutions can create a useful reuse system while weakening a claim of novel-task reasoning. Separate the producer and consumer splits, record artifact timestamps, and audit duplicate or templated trajectories. If synthetic queries are generated from known evidence, evaluate whether their wording makes the source unusually easy to retrieve.

Longitudinal stores introduce order dependence. A task can become easier because a previous evaluation task inserted its answer into memory. This may be an intended continual-learning setting, but it is not equivalent to an independent held-out test. Declare whether updates are allowed, reset stores between conditions where required, and report frozen-index and evolving-index results separately. Paired direct and indirect questions should read the same frozen state when the goal is to diagnose implicit relevance.

Uncertainty should include tasks, corpus samples, and stochastic access policies. Repeating decoding on a single fixed corpus does not estimate robustness to document composition. Rebuilding an index with a different embedding model can change candidate sets and calibration, not just vector quality. Treat index generation as part of the configuration. Judge-based relevance and citation scoring should be checked against an independently reviewed subset, with disagreement retained.

None of the proposed evaluations has been run for this monograph. The reviewed 2026 methods provide concrete mechanisms and author-reported evidence under their stated settings. A rigorous survey can compare those mechanisms, identify causal confounds, and derive discriminating tests without presenting new measurements. That boundary is especially important when active retrieval is advertised as a general solution to memory: its benefits remain conditional on what was stored, how it was indexed, and what the policy can infer.

